%%%%%%%%%%%%%%%%%%%%%%%%%%%%%%%%%%%%%%%%%%%%%%%%%%
\needspace{12\baselineskip}
\section{Formale Sprachen und Grammatiken}\label{sec:grammatiken}
%%%%%%%%%%%%%%%%%%%%%%%%%%%%%%%%%%%%%%%%%%%%%%%%%%
\worum
Natürliche Sprache ist ungenau und mehrdeutig (\enquote{Rasen ist gefährlich!}). Eine Programmiersprache muss dagegen eindeutig festlegen, welche Texte gültige Programme sind (\textbf{Syntax}) und was sie bedeuten (\textbf{Semantik}). Die Syntax beschreibt man mit einer \textbf{Grammatik}: einer Menge von Ersetzungsregeln, mit denen sich aus einem Startsymbol alle gültigen Wörter der Sprache erzeugen lassen. Für dieselbe Sache gibt es drei gleich mächtige Schreibweisen: die formale Grammatik nach Chomsky, die \textbf{EBNF} und das \textbf{Syntaxdiagramm}.

\begin{merke}[Grundbegriffe]
\begin{itemize}[itemsep=0.1em]
\item \textbf{Alphabet} $A$: endliche Menge von Symbolen. \ \textbf{Wort}: endliche Folge von Symbolen aus $A$. \\
$\varepsilon$ ist das leere Wort, $A^*$ die Menge aller Wörter über $A$. Eine \textbf{Sprache} ist eine Teilmenge $L \subseteq A^*$.
\item \textbf{Grammatik} $G = (N, T, P, S)$ mit
	\begin{itemize}[itemsep=0em]
	\item $N$: Nichtterminalsymbole (werden noch ersetzt), \ $T$: Terminalsymbole (bleiben stehen), $N \cap T = \emptyset$
	\item $P$: Produktionen (Regeln) der Form $(p, q)$, gelesen \enquote{ersetze $p$ durch $q$}
	\item $S \in N$: Startsymbol
	\end{itemize}
\item $L(G)$ ist die Menge aller Wörter, die nur aus Terminalsymbolen bestehen und sich aus $S$ ableiten lassen.
\end{itemize}
\end{merke}

\begin{merke}[EBNF]
\begin{tabular}{@{}ll@{\qquad}ll@{}}
\code{"a"} & Terminalsymbol & \code{[ A ]} & Option: A darf einmal oder gar nicht vorkommen \\
\code{A B} & Folge: erst A, dann B & \code{\{ A \}} & Wiederholung: A beliebig oft, auch nullmal \\
\code{A | B} & Alternative: A oder B & \code{( A )} & Gruppierung
\end{tabular}
\end{merke}

\begin{merke}[Chomsky-Hierarchie]
Je höher der Typ, desto stärker sind die Regeln eingeschränkt und desto einfacher lässt sich prüfen, ob ein Wort zur Sprache gehört.

\smallskip
\begin{tabular}{@{}lL{3.3cm}L{9.2cm}@{}}
\textbf{Typ 0} & allgemein & beliebige Regeln \\
\textbf{Typ 1} & kontextsensitiv & Regeln verkürzen nicht: $|p| \leq |q|$. Auf der linken Seite dürfen neben dem Nichtterminal weitere Zeichen (der \enquote{Kontext}) stehen. \\
\textbf{Typ 2} & kontextfrei & links steht genau ein Nichtterminal \\
\textbf{Typ 3} & regulär & zusätzlich: rechts steht höchstens ein Nichtterminal, und zwar immer am Ende (rechtslinear) oder immer am Anfang (linkslinear)
\end{tabular}

\smallskip
Programmiersprachen sind in der Regel vom Typ 2 oder Typ 3. EBNF und Syntaxdiagramme beschreiben genau die kontextfreien Sprachen.
\end{merke}

\beispiele
\beispiel{Wörter ableiten und die Sprache erkennen}
Gegeben sei $G$ mit $N = \{S, A\}$, $T = \{x, y\}$ und $P = \{(S, xA),\ (A, yS),\ (A, y)\}$. Man beginnt beim Startsymbol und wendet Regeln an, bis kein Nichtterminal mehr übrig ist:
\[
S \to xA \to xy
\qquad\qquad
S \to xA \to xyS \to xyxA \to xyxy
\]
An den Beispielen erkennt man: $L(G) = \{\, w \mid w = (xy)^n,\ n \in \mathbb{N} \,\}$. Links steht immer genau ein Nichtterminal, rechts steht das Nichtterminal immer am Ende: Die Grammatik ist vom Typ 3.

\beispiel{Zu einer Sprache die Grammatik finden}
Gesucht ist $P$ für $L = \{\, w \mid w = a^n b^{2n},\ n \in \mathbb{N} \,\}$, also $abb$, $aabbbb$, \ldots

Die Anzahl der $a$ und $b$ hängt zusammen. Solche Wörter baut man \enquote{von außen nach innen} auf: Jede Anwendung der Regel erzeugt vorne ein $a$ und hinten zwei $b$.
\[
P = \{(S, aSbb),\ (S, abb)\}
\qquad\text{z.\,B.}\quad S \to aSbb \to aabbbb
\]
Soll auch $n = 0$ erlaubt sein ($n \in \mathbb{N}_0$), muss das leere Wort erzeugt werden können: $P = \{(S, aSbb),\ (S, \varepsilon)\}$.

Steht vor oder hinter dem wiederholten Teil noch ein fester Teil, lagert man die Wiederholung in ein eigenes Nichtterminal aus. Für $L = \{\, w \mid w = c\,a^n b^n\,d,\ n \in \mathbb{N} \,\}$:
\[
N = \{S, B\}, \qquad P = \{(S, cBd),\ (B, aBb),\ (B, ab)\}
\]

\beispiel{EBNF prüfen}
Können die Wörter mit den EBNF-Ausdrücken erzeugt werden?
\begin{center}
\renewcommand{\arraystretch}{1.3}
\begin{tabular}{|l|C{2.6cm}|C{2.6cm}|C{2.6cm}|}
\hline
 & \code{\{a\} b} & \code{a [b] \{a\}} & \code{\{ab | b\}} \\ \hline
\code{ab} & ja & ja & ja \\ \hline
\code{aab} & ja & nein & nein \\ \hline
\code{bab} & nein & nein & ja \\ \hline
$\varepsilon$ & nein & nein & ja \\ \hline
\end{tabular}
\end{center}
Zur Begründung zerlegt man das Wort passend zum Ausdruck. \code{aab} passt zu \code{\{a\} b}: zweimal \code{a}, dann \code{b}. Zu \code{a [b] \{a\}} passt es nicht, denn nach dem optionalen \code{b} dürfen nur noch \code{a} folgen. In \code{\{ab | b\}} muss auf jedes \code{a} direkt ein \code{b} folgen. Das leere Wort entsteht, wenn eine Wiederholung nullmal durchlaufen wird; das geht nur, wenn der Ausdruck kein Pflichtzeichen enthält.

\beispiel{Vom Syntaxdiagramm zu Grammatik und EBNF}
Ein Wort entsteht, indem man das Diagramm links betritt, den Pfeilen folgt, die Symbole der durchlaufenen Kreise notiert und das Diagramm rechts verlässt. (Kreise stehen für Terminalsymbole, Rechtecke für Nichtterminale.)

\noindent\begin{minipage}[c]{0.42\linewidth}
\begin{tikzpicture}[syntax]
	\node[knoten] (a) at (0,0) {\itshape a};
	\node[knoten] (b) at (1.7,0) {\itshape b};
	\node[knoten] (c) at (3.4,0) {\itshape c};
	\draw[linie, ->] (-1.1,0) -- (a);
	\draw[linie, ->] (a) -- (b);
	\draw[linie, ->] (b) -- (c);
	\draw[linie, ->] (c) -- (4.6,0);
	\draw[linie, ->] (c.south) -- (3.4,-0.9) -- (1.7,-0.9) -- (b.south);
\end{tikzpicture}
\end{minipage}\hfill
\begin{minipage}[c]{0.55\linewidth}
Mögliche Wörter: $abc$, $abcbc$, $abcbcbc$, \ldots

\smallskip
\textbf{EBNF:} \quad \code{"a" "b" "c" \{ "b" "c" \}}

\smallskip
\textbf{Grammatik:} \ $N = \{S, B\}$, \ $T = \{a, b, c\}$, \ Startsymbol $S$, \\
$P = \{(S, aB),\ (B, bcB),\ (B, bc)\}$
\end{minipage}

\aufgaben
\aufgabe{Sprache bestimmen}
Bestimmen Sie jeweils (mithilfe von Beispielableitungen) die Sprache $L(G)$ und geben Sie den Chomsky-Typ der Grammatik an. Es gilt $T = \{a, b\}$.
\begin{enumerate}[label=\alph*)]
\item $N = \{S\}$, \ $P = \{(S, b),\ (S, aSa)\}$
\item $N = \{S, B\}$, \ $P = \{(S, aB),\ (B, bS),\ (B, b)\}$
\item $N = \{S, B\}$, \ $P = \{(S, aaS),\ (S, B),\ (B, bB),\ (B, b)\}$
\end{enumerate}

\aufgabe{Grammatik bestimmen}
Sei $G_i = (N, T, P_i, S)$ mit $T = \{a, b, c, x\}$. Bestimmen Sie $N$ und $P_i$ möglichst minimal so, dass gilt:
\begin{enumerate}[label=\alph*)]
\item $L(G_1) = \{\, w \mid w = a^{n} b^{3n},\ n \in \mathbb{N} \,\}$
\item $L(G_2) = \{\, w \mid w = (ba)^{n} b,\ n \in \mathbb{N}_0 \,\}$
\item $L(G_3) = \{\, w \mid w = a^{n} (cb)^{2n} xx,\ n \in \mathbb{N} \,\}$
\item $L(G_4) = \{\, w \mid w = (ccc)^{n} b^3 x^n,\ n \in \mathbb{N}_0 \,\}$
\end{enumerate}

\aufgabe[\K]{Grammatik bestimmen}
Gegeben sei die Grammatik $G = (N, T, P, S)$ mit $T = \{a, b, c\}$. Bestimmen Sie $N$ und $P$ möglichst minimal so, dass die folgende Sprache erzeugt wird:
\[
L(G) = \{\, w \mid w = a\,b^{3n} c^{n},\ n \in \mathbb{N}_0 \,\}
\]

\aufgabe[\K]{EBNF prüfen}
Überprüfen Sie, ob die folgenden Wörter mit den vorgegebenen EBNF-Ausdrücken erzeugt werden können, und tragen Sie \emph{ja} oder \emph{nein} ein. (Leerzeichen dienen nur der Lesbarkeit.)
\begin{center}
\renewcommand{\arraystretch}{1.5}
\begin{tabular}{|l|C{2.8cm}|C{2.8cm}|C{2.8cm}|}
\hline
 & \code{\{a | ba\}} & \code{a \{b\} [a]} & \code{\{a [b]\} b} \\ \hline
\code{ab} & & & \\ \hline
\code{aba} & & & \\ \hline
\code{baab} & & & \\ \hline
\code{abb} & & & \\ \hline
$\varepsilon$ & & & \\ \hline
\end{tabular}
\end{center}

\aufgabe{EBNF mit mehreren Regeln}
Gegeben seien die folgenden Regeln in EBNF:
\begin{ebnf}
	Ziffer & \q{0} | \q{1} | \ldots\ | \q{9} \\
	Buchstabe & \q{a} | \q{b} | \ldots\ | \q{z} \\[2mm]
	Start & Kennung | Preis \\
	Kennung & Buchstabe \{ Buchstabe \} \q{-} Ziffer Ziffer [ Ziffer ] \\
	Preis & Ziffer \{ Ziffer \} [ \q{,} Ziffer Ziffer ] \q{€}
\end{ebnf}
Lassen sich die folgenden Wörter darstellen? Begründen Sie Ihre Antwort.
\begin{multicols}{5}
\begin{enumerate}[label=\alph*)]
\item \code{abc-12}
\item \code{x-1234}
\item \code{12,5€}
\item \code{7€}
\item \code{-12}
\end{enumerate}
\end{multicols}

\aufgabe{EBNF entwerfen}
Definieren Sie jeweils eine Menge von EBNF-Regeln.
\begin{enumerate}[label=\alph*)]
\item ganze Zahlen mit optionalem Minuszeichen und ohne führende Nullen (\code{0}, \code{-7}, \code{120}, aber nicht \code{007})
\item gerade Dualzahlen ohne führende Nullen; die kleinste Zahl ist 2 (\code{10}, \code{100}, \code{110}, \ldots)
\item gültige Kalendertage der Form \textbf{TT.MM.} von \code{01.01.} bis \code{31.12.} (Die unterschiedliche Länge der Monate dürfen Sie ignorieren.)
\end{enumerate}

\aufgabe[\K]{Syntaxdiagramm}
Das folgende Syntaxdiagramm beschreibt eine Sprache $L$.
\begin{center}
\begin{tikzpicture}[syntax]
	\node[knoten] (x1) at (0,0) {0};
	\node[knoten] (x2) at (2.2,0) {1};
	\node[knoten] (x4) at (4.4,0) {1};
	\node[knoten] (x3) at (2.2,-1.8) {0};
	\draw[linie, ->] (-1.3,0) -- (x1);
	\draw[linie, ->] (x1) -- (x2);
	\draw[linie, ->] (x2) -- (x4);
	\draw[linie, ->] (x4) -- (5.7,0);
	\draw[linie, ->] (x2) to[bend right=35] (x3);
	\draw[linie, ->] (x3) to[bend right=35] (x2);
	\draw[linie, ->] (x3) to[out=240, in=300, looseness=7] (x3);
\end{tikzpicture}
\end{center}
\begin{enumerate}[label=\alph*)]
\item Geben Sie drei Wörter der Sprache an.
\item Geben Sie zu dieser Sprache eine Grammatik $G = (N, T, P, S)$ an. Definieren Sie $N$, $T$, $P$ und $S$.
\item Geben Sie zu der Sprache auch einen EBNF-Ausdruck an.
\end{enumerate}

\aufgabe{Syntaxdiagramm zeichnen}
Zeichnen Sie zu den folgenden EBNF-Ausdrücken jeweils ein Syntaxdiagramm.
\begin{multicols}{3}
\begin{enumerate}[label=\alph*)]
\item \code{"a" \{ "b" | "c" \} "d"}
\item \code{[ "-" ] "1" \{ "0" \}}
\item \code{"x" \{ "y" "x" \} [ "z" ]}
\end{enumerate}
\end{multicols}

\aufgabe[\K]{Chomsky-Hierarchie}
\begin{enumerate}[label=\alph*)]
\item Was versteht man unter einer kontextsensitiven Grammatik? Worin unterscheidet sie sich von einer kontextfreien Grammatik?
\item Ordnen Sie die folgenden Produktionsmengen dem höchstmöglichen Chomsky-Typ zu. \\
$P_1 = \{(S, aSb),\ (S, ab)\}$ \qquad $P_2 = \{(S, aS),\ (S, b)\}$ \\
$P_3 = \{(S, aSBc),\ (S, abc),\ (cB, Bc),\ (bB, bb)\}$ \qquad $P_4 = \{(S, aSb),\ (aSb, a)\}$
\end{enumerate}

\aufgabe[\S]{Mehr als kontextfrei}
Die Produktionen $P_3$ aus der vorigen Aufgabe erzeugen die Sprache $\{\, a^n b^n c^n \mid n \in \mathbb{N} \,\}$, die sich mit keiner kontextfreien Grammatik beschreiben lässt. Leiten Sie mit $P_3$ das Wort $aabbcc$ ab und geben Sie bei jedem Schritt die verwendete Regel an.

\ueben
\aufgabe{Ableiten}
\begin{enumerate}[label=\alph*)]
\item Gegeben sei $P = \{(S, aSb),\ (S, ab)\}$. Leiten Sie das Wort $aaabbb$ ab.
\item Gegeben sei $P = \{(S, AB),\ (A, aA),\ (A, a),\ (B, bB),\ (B, b)\}$. Leiten Sie die Wörter $aab$ und $abbb$ ab und beschreiben Sie die Sprache.
\end{enumerate}

\aufgabe{Gehört das Wort zur Sprache?}
Gegeben sei $L = \{\, w \mid w = a^{n} b^{2n} c,\ n \in \mathbb{N}_0 \,\}$. Welche der folgenden Wörter gehören zu $L$?
\begin{multicols}{5}
\begin{enumerate}[label=\alph*)]
\item \code{c}
\item \code{abbc}
\item \code{aabbc}
\item \code{aabbbbc}
\item \code{abb}
\end{enumerate}
\end{multicols}

\aufgabe{Grammatiken finden}
Geben Sie jeweils $N$ und $P$ an ($T$ ergibt sich aus der Sprache).
\begin{enumerate}[label=\alph*)]
\item $L_1 = \{\, a^n b^n \mid n \in \mathbb{N} \,\}$
\item $L_2 = \{\, a^n b^m \mid n, m \in \mathbb{N} \,\}$
\item $L_3 = \{\, (ab)^n c^n \mid n \in \mathbb{N}_0 \,\}$
\item $L_4 = \{\, x\,a^{2n}\,y \mid n \in \mathbb{N}_0 \,\}$
\item $L_5 = \{\, a^n\,b\,a^n \mid n \in \mathbb{N}_0 \,\}$
\end{enumerate}

\aufgabe{EBNF prüfen}
\begin{center}
\renewcommand{\arraystretch}{1.4}
\begin{tabular}{|l|C{2.8cm}|C{2.8cm}|C{2.8cm}|}
\hline
 & \code{[a] \{b\}} & \code{\{a b\} a} & \code{(a | b) \{a\} [b]} \\ \hline
\code{a} & & & \\ \hline
\code{abb} & & & \\ \hline
\code{aba} & & & \\ \hline
\code{baab} & & & \\ \hline
$\varepsilon$ & & & \\ \hline
\end{tabular}
\end{center}

\aufgabe{EBNF schreiben}
Definieren Sie jeweils EBNF-Regeln.
\begin{enumerate}[label=\alph*)]
\item eine deutsche Postleitzahl (genau fünf Ziffern)
\item eine Dezimalzahl mit Komma und mindestens einer Stelle davor und danach (\code{3,14}, \code{0,5})
\item ein vereinfachtes Kfz-Kennzeichen: ein bis drei Großbuchstaben, ein Bindestrich, ein bis zwei Großbuchstaben, ein Leerzeichen und eine Zahl mit ein bis vier Ziffern, die nicht mit 0 beginnt
\end{enumerate}

\aufgabe{Vom Diagramm zur EBNF}
Geben Sie zu dem Syntaxdiagramm drei Wörter der Sprache, einen EBNF-Ausdruck und eine Grammatik an.
\begin{center}
\begin{tikzpicture}[syntax]
	\node[knoten] (a) at (0,0) {\itshape a};
	\node[knoten] (b) at (0,-1.5) {\itshape b};
	\draw[linie, ->] (-1.4,0) -- (a);
	\draw[linie, ->] (a) -- (1.6,0);
	\draw[linie, ->] (a) to[bend left=40] (b);
	\draw[linie, ->] (b) to[bend left=40] (a);
\end{tikzpicture}
\end{center}

\aufgabe{Chomsky-Typ bestimmen}
Geben Sie zu jeder Produktionsmenge den höchstmöglichen Typ an.
\begin{multicols}{2}
\begin{enumerate}[label=\alph*)]
\item $\{(S, aS),\ (S, a)\}$
\item $\{(S, aSa),\ (S, bSb),\ (S, c)\}$
\item $\{(S, Sa),\ (S, b)\}$
\item $\{(S, aAb),\ (aA, aaA),\ (A, b)\}$
\end{enumerate}
\end{multicols}
